\chapter*{后记}
\addcontentsline{toc}{chapter}{后记} % 手动加入目录
\markboth{后记}{后记} % 设置页眉

2002年，一个很偶然的机会，我对中国苏维埃运动的历史产生了兴趣，并在此之后不无后见之明地发现，因为出生在江西赣州的关系，使我对这段历史的研究多了些不可选择的天然条件。自幼的耳濡目染，可以让我更近、更直观、更便捷地体会到革命的原初形态，童年时代农村生活的记忆，中学时候学校后山红军挖掘的蜿蜒的坑道，大学期间读到的关于赣南农村社会状况的国外研究论文，似乎都一起复活，指引着我去追寻当年那场曾在故乡激起山乡巨变的熊熊烈火。


2003年，申请关于第五次“围剿”和反“围剿”的研究课题获准，到2005年基本写出初稿。此后，就是没完没了地补充和修改过程。没有办法知道存了多少稿，只知道从开始写作到现在，日常处理文字的电脑已经换到了第四台。


从2002年到现在，正好是十年。接下来，好像就该说“十年磨一剑”了。不过，不想也不能这么说。因为这十年，当然不仅仅是为这本书而存在，甚至还有几年，连学术也被悄悄地放到了一边。生活就像历史一样，总是充满着困惑和选择，只是这些，仅对于个体的自我有意义罢了。


感谢那些始终予我以鼓励和帮助的人们，没有他们，这本书不知道能在什么时候完成，甚至我的研究也不知道还会不会继续下去。对于一个常常是处于被动状态的消极者而言，有时候，一句简单的肯定、轻声的慰藉，也许都会有令人难以忘怀的效果。


对于这些无私的帮助，我从不知道何以为报。更主动地投入生活，更潜心地研究学术，更友善地面对他人，或许就是他们期望得到的回应吧。


感谢父母的养育、妻子的支持、儿子的进步。感谢阅读在给了我一副眼镜的同时，还给了一双我自己的眼睛。

\begin{flushright}
2011年初春
\end{flushright}